\documentclass[12pt]{article}

\usepackage{fullpage,amsmath,amssymb}

\begin{document}
\begin{center}
\textbf{Mars Hill University}
\\
\textbf{Math 107: Finite Mathematics, Section 01 (3 credits)}
\\Fall 2013
\end{center}

\noindent\textbf{Instructor Information:}

\indent \textit{Instructor:} Dr. Laura Steil, Ph.D.
\\ \indent \textit{Office:} Wall 408
\\ \indent \textit{Phone:} 689-1168
$*$ Please note that according to college policy, instructors cannot return phone calls to students with long-distance phone numbers (i.e., cell phones).  If you cannot leave a local phone number, the instructor will respond to your phone message via email.
\\ \indent \textit{Email:} lsteil@mhu.edu
\\ \indent \textit{Office Hours:}
\\ \indent Monday 11:00am-noon, 2:00-4:00pm
\\ \indent Wednesday 11:00am-noon, 2:00-4:00pm
\\ \indent Thursday 1:00-2:00pm
\\ \indent Friday 11:00am-noon
\\ \indent Other times available by appointment

\vspace{.1 in}
\noindent\textbf{Time and Location:} MWF 9:00-9:50am,  Wall 422
\vspace{.1 in}

 \noindent\textbf{Required Materials:}
\\ \indent\textbf{Textbook:}
\textit{ Mathematics, A practical Odyssey, 7th edition, } Johnson, David B. and Thomas A. Mowry
   \textit{ISBN}: 0840049137
\noindent
\\ \indent\textbf{Calculator: }
      You will need a scientific calculator for this course.
\\ \indent\textbf{WebAssign:} Most homework assignments will be turned in through the online homework system WebAssign.  You will need to purchase access to this system.  More information will be given on a handout.
\vspace{.1 in}

 \noindent\textbf{Course Description:} Math 107 is an introduction to the following four related subjects: combinatorics, symbolic logic, probability theory, and descriptive statistics.
\vspace{.1 in}

 \noindent\textbf{Course Goals:}  Students who complete Math 107 will have:
      \begin{itemize}
\item An ability to use counting formulas, permutation numbers, and combination numbers in solving combinatoric problems.
    \item  An understanding of basic symbolic logic.
\item Skill in using Venn diagrams and truth tables to test the validity of arguments.
\item A knowledge of basic probability theory and skill in applying probability formulas in various situations.
\item An understanding of basic descriptive statistics including the ability to compute means, medians, and standard deviations, and an ability to construct histograms.
\item A knowledge of the normal distribution and the ability to apply it.
    \end{itemize}
\vspace{.1 in}

\noindent
\textbf{Grading:}
Course grades are dependent upon class tests, homework and quizzes, and a final examination. There will be four in-class tests during the semester.  Each of these tests is worth 100 points.  At the end of the semester, a comprehensive final examination will be given.  The exam is worth 200 points.  Homework cannot be made up.  There will be no make-up tests unless you are participating in an official school function.  In this case, you should notify the instructor in advance to make the appropriate arrangements.  Otherwise, if you miss one chapter test, your exam grade will count for that test.  If you miss more than one chapter test, you will receive a zero as your test grade on each additional test.



\hspace{1in}
\begin{tabular}[t]{llll}
Homework and Quizzes &&& 100 points \\
4 class tests & & & 400 points \\
Final Exam  & & & 200 points \\
\end{tabular}

\noindent
Your grade will be based on the following point system:

\hspace{1in}
\begin{tabular}[t]{ll}
A+ & 686-700 points \\
A & 651-685 points \\
A- & 630-650 points\\
B+ & 616-629 points\\
B & 581-615 points \\
B- & 560-580 points \\
C+ & 549-559 points \\
C & 511-548 points \\
C- & 490-510 points \\
F & 0-489 points
\end{tabular}

\textit{Snow Policy:}
If class is canceled due to inclement weather, I will send an email to the class via MHU Webmail.

\textit{It is your responsibility to check your MHU email account regularly for announcements from the instructor.}

\vspace{.1 in}
\noindent\textbf{Homework and Quizzes:}
You will have online homework assignments through WebAssign.  There will be occasional written homework assignments to turn in as well.  You will need access to WebAssign for this course.  See the handout about logging onto WebAssign to start.  Each homework assignment has a posted due date and time.  Late submissions or make up work are not allowed.  Quizzes will be given during the semester at the discretion of the instructor.   At the end of the semester at least two assignment grades will be dropped to allow for unavoidable situations.  Then the individual homework assignment and quiz grades will be averaged together (with each assignment given equal weight) to your score out of 100 points.

\vspace{.1 in}
\noindent\textbf{Tests and Final Exam:}
 We will have four in class tests and one
final exam.  The
final exam will be given on ({\bf Friday, December 13 at 8am}). The final exam is comprehensive.

\vspace{.1 in}
 \noindent\textbf{Attendance Policy:}
A student who has been absent more than five times during the semester will have 15 points of the possible 700 points subtracted for each absence over five.  Three tardies will be counted as one absence.  Early departure from class counts as a tardy, as well as any departure from class (even if you do return).  \underline{There are no excused absences.}  Absences for illness, athletic events, band and choir tours, medical appointments, weddings, funerals, etc., are included in the five absences which are not penalized.  If you are an athlete and know that you will miss more than five classes due to your game schedule, please talk to me ASAP.

\vspace{.1 in}
\noindent\textbf{Disability Accommodations:} If you have a disability that requires accommodations in this course, it is your responsibility to file the appropriate paperwork with the Disabilities Office and/or Student Support Services.  Once the Disabilities Office or Student Support Services has approved the appropriate accommodations, it is again your responsibility to bring this documentation to me at least one week prior to receiving accommodations.

\vspace{.1 in}
\noindent\textbf{Classroom Behavior:}
You are expected to arrive on time for class and stay until the class is dismissed.  Please turn off your cell phone and put it away before class begins. Texting during class may result in your being asked to leave. Please do not bring food into the building, and only beverages with a lid will be allowed. No tobacco products of any kind may be used inside of this building.

\vspace{.1 in}
\noindent\textbf{Honor Code:}
 All Mars Hill University students are expected to conform to the following Honor Code:

\begin{center}Honor Code\end{center}
	We, the students of Mars Hill University, pledge ourselves to uphold integrity, honesty,
and academic responsibility in and out of the classroom.

\begin{center}Honor Pledge\end{center}
	On my honor, I have neither given nor received any academic aid or information
	that would violate the Honor Code of Mars Hill University.

\newpage

\thispagestyle{empty}
\begin{center}
\begin{tabular}{|c|c|l|l|}
\hline
&\textbf{Tentative Schedule}&\\
 \hline
Wednesday & August 28 & Syllabus, Intro. to WebAssign, Section 2.1\\

\hline
Friday & August 30 & Section 2.1\\

\hline
Monday & September 2 & Section 2.2\\

\hline
Wednesday & September 4 &Section 2.3\\

\hline
Friday & September 6 & Section 2.3\\

 \hline
Monday & September 9& Section 2.4\\

 \hline
Wednesday & September 11& Section 2.4\\

 \hline
Friday &September 13 & Review \\

 \hline
Monday & September 16& \textbf{Test 1}\\

 \hline
Wednesday & September 18& Section 1.1\\

 \hline
Friday &September 20 & Section 1.2 \\

 \hline
Monday & September 23& Section 1.2\\

 \hline
Wednesday & September 25& Section 1.3\\

 \hline
Friday &September 27 & Section 1.3 \\

 \hline
Monday & September 30& Section 1.4\\

 \hline
Wednesday & October 2& Section 1.5\\

 \hline
Friday & October 4 & Section 1.5 \\

 \hline
Monday & October 7 & Review\\

\hline
Wednesday & October 9 & \textbf{Test 2}\\

\hline
Friday & October 11 & Section 3.2\\

\hline
Monday & October 14 & \textbf{Fall Break}\\

\hline
Wednesday & October 16 & Section 3.3\\

\hline
Friday & October 18 & Section 3.3 \\

\hline
Monday & October 21 & Section 3.4\\

\hline
Wednesday & October 23 & Section 3.4\\

\hline
Friday & October 25 & Section 3.5 \\

\hline
Monday & October 28 & Section 3.6\\

\hline
Wednesday & October 30 & Section 3.7\\

\hline
Friday & November 1 & Review \\

\hline
Monday & November 4 & \textbf{Test 3}\\

\hline
Wednesday & November 6 & Section 4.1\\

\hline
Friday & November 8 & Section 4.2 \\

\hline
Monday & November 11 & Section 4.2-4.3 \\

\hline
Wednesday & November 13 & Section 4.3\\

\hline
Friday & November 15 & Section 4.3-4.4 \\

\hline
Monday & November 18 & Section 4.4 \\

\hline
Wednesday & November 20 & Section 4.4\\

\hline
Friday & November 22 & Section 4.4\\

\hline
Monday & November 25 & Section 4.4\\

\hline
Wednesday & November 27 & \textbf{Thanksgiving Break}\\

\hline
Friday & November 29 & \textbf{Thanksgiving Break}\\

\hline
Monday & December 2 & Review\\

\hline
Wednesday & December 4 & \textbf{Test 4} \\

\hline
Friday & December 6 & Review\\

\hline
Monday & December 9 & Review\\
 \hline

\end{tabular}
\end{center}

\end{document} 